% Ad Atticum 3.15 --- headnote
% ad-atticum-03-15, 19 August 58 BC, written at Thessalonica

\textit{Cicero to Atticus, written from Thessalonica on the
fourteenth day before the Kalends of September (19 August)
58 BC --- the longest of the surviving Thessalonican letters
and the bitterest. Four letters from Atticus have arrived
together on the Ides of August, and Cicero answers each in
turn: \textsection 2 to the rebuke for weakness of mind,
\textsection 3 to the news of the Senate proceedings (Curio
read out a hostile speech --- whose origin is unknown ---
which contradicts what Atticus had written), \textsection 4
to Varro's report of Pompey, opening into Cicero's longest
self-rebuke and most pointed reproach to Atticus.}

\textit{The famous lines are in \textsection 4--5: ``\textit{caeci,
caeci inquam fuimus}'' --- ``Blind, blind I say we were'' in
the public mourning and the supplication of the people; and
the regret that no one was at hand, when he was terrified by
Pompey's ``stinted answer,'' to call him back from a
disgraceful counsel --- ``either I should have fallen with
honour, or we should be living today as victors.''
\textsection 7 turns directly on Atticus: ``you, who, even if
you did not have more talent than I had, were certainly less
frightened.'' Atticus, the letter implies, is now bearing in
the shipwreck the labours that earlier presence of mind would
have spared him; the warmth of \textsection 8 takes the edge
off but does not retract the charge. The technical part of
the letter is \textsection 5--6 on the Clodian
\textit{privilegium} --- the law against Cicero by name ---
and whether to abrogate it directly or set it aside through
a senatorial decree.}
